\section{Dilation, distribution traces and exact logarithmic contacts}
\label{dilfp:sec:main}
The contact theorem answers an extension question suggested by the
correlation reports: classical distribution rows need a normalization
correction when interpreted as finite-part identities. Fourier traces
make the correction computable at every Stieltjes and derivative order.
This also describes how dilation acts on the one-generator convolution
algebra, without any arithmetic-independence assumption.

For a positive integer $q$, define the trace $\mathcal S_q$ and covering
pullback $\mathcal P_q$ on periodic distributions by
\[
 \widehat{\mathcal S_qT}(k)=q\widehat T(qk),\qquad
 \widehat{\mathcal P_qT}(k)=
 \begin{cases}\widehat T(k/q),&q\mid k,\\0,&q\nmid k.\end{cases}
\]
For ordinary periodic functions these are
$\mathcal S_qf(x)=\sum_{j=0}^{q-1}f((x+j)/q)$ and
$\mathcal P_qf(x)=f(\{qx\})$. Put $\ell_q=\log q$ and
$\mathcal L_q=q^{-1}\mathcal S_q$.

\begin{theorem}[The trace translates the convolution generator]
\label{dilfp:thm:trace}
For every $m\ge0$,
\begin{equation}\label{dilfp:eq:trace}
 \mathcal S_qG_m
 =q\sum_{j=0}^m\binom mj(-\ell_q)^{m-j}G_j
   +\frac{q(-\ell_q)^{m+1}}{m+1}\stconvId.
\end{equation}
The normalized trace $\mathcal L_q$ is a unital homomorphism of the
mean-zero convolution algebra. On its generator $X=G_0$ it acts by
$X\mapsto X-\ell_q\stconvId$. These actions compose as
$\mathcal L_q\mathcal L_r=\mathcal L_{qr}$.
\end{theorem}
\begin{proof}
The Fourier definition gives
$\mathcal S_qR_\alpha=q^{1-\alpha}R_\alpha$. Apply this to the
Laurent family \eqref{stconv:eq:ZLaurent}:
\[
 \mathcal S_q\left(\frac{\stconvId}{z}
                  +\sum_{m\ge0}G_m\frac{z^m}{m!}\right)
 =q e^{-\ell_qz}\left(\frac{\stconvId}{z}
                  +\sum_{m\ge0}G_m\frac{z^m}{m!}\right).
\]
The regular coefficients give \eqref{dilfp:eq:trace}, including the
term produced by the pole. On Fourier coefficients,
$\mathcal L_q(T*S)=(\mathcal L_qT)*(\mathcal L_qS)$ and
$\mathcal L_q\stconvId=\stconvId$. The generator formula is the
$m=0$ case, and repeated traces evaluate a Fourier coefficient at $qrk$.
\end{proof}

\begin{theorem}[Finite part does not commute with covering pullback]
\label{dilfp:thm:cover}
Let $\operatorname{Fp}_x\gamma_m(\{qx\})$ use the constant term in
the unscaled local coordinate $x-j/q$ at each of the $q$ singularities,
with right-hand cutoffs tending independently to zero. Then
\begin{equation}\label{dilfp:eq:cover}
 \mathcal P_qG_m
 =\operatorname{Fp}_x\gamma_m(\{qx\})
 +\frac{\ell_q^{m+1}}{q(m+1)}
                    \sum_{j=0}^{q-1}\delta_{j/q}.
\end{equation}
In particular the raw-coordinate finite part has mean
$-\ell_q^{m+1}/(m+1)$, while the covering pullback has mean zero.
\end{theorem}
\begin{proof}
On the $j$th segment substitute $y=q(x-j/q)$. Pulling back the
cutoff prescription \eqref{stconv:eq:Gcutoff} gives its local
counterterm
\[
 \frac{\log^{m+1}(q\epsilon)}{q(m+1)}\phi(j/q).
\]
The raw-coordinate constant-term prescription keeps only the divergent
powers of $\log\epsilon$ in this expression. Their difference is its
constant coefficient, $\ell_q^{m+1}\phi(j/q)/(q(m+1))$.
All Taylor remainder integrals converge as in the definition of $G_m$.
Summing the segments proves \eqref{dilfp:eq:cover}. The Fourier
pullback preserves the zeroth coefficient, and pairing the contact sum
with the constant function proves the mean statement.
\end{proof}

For the trace itself, put $c_{m,q}=-q(-\ell_q)^{m+1}/(m+1)$.
The classical pointwise multiplication row and its raw-coordinate
extension read
\[
 \operatorname{Fp}_x[\mathcal S_q\gamma_m]
 =q\sum_{j=0}^m\binom mj(-\ell_q)^{m-j}G_j+c_{m,q}\,1.
\]
Comparing with \eqref{dilfp:eq:trace} gives the separate identity
\begin{equation}\label{dilfp:eq:trace-commutator}
 \mathcal S_qG_m
 =\operatorname{Fp}_x[\mathcal S_q\gamma_m]-c_{m,q}\delta_0.
\end{equation}
The logarithmic term thus has different signs and support for covering
pullback and trace; substituting one convention for the other is invalid.

\begin{theorem}[Every differentiated trace and its contact]
\label{dilfp:thm:derivative}
Use $F_{m,r}$ and $a_{m,r}=(-1)^{m+1}m!e_{m+1}(r)$ from the
contact theorem, and let $r\ge1$. Then
\begin{align}\label{dilfp:eq:derivative}
 \mathcal S_qF_{m,r}
 ={}&q^{r+1}\sum_{j=0}^m\binom mj(-\ell_q)^{m-j}F_{j,r}
                         +q^{r+1}\kappa_{m,r}(\ell_q)\delta_0^{(r)},\notag\\
 \kappa_{m,r}(\ell)
 ={}&\frac{(-\ell)^{m+1}}{m+1}
       +\sum_{j=0}^m\binom mj(-\ell)^{m-j}a_{j,r}-a_{m,r}.
\end{align}
The first sum is exactly
$\operatorname{Fp}_x[\mathcal S_q\gamma_m^{(r)}]$.
For the canonical polygamma finite parts, the particularly simple
all-order specialization is
\begin{equation}\label{dilfp:eq:polygamma}
 \mathcal S_q\Psi_r=q^{r+1}
                   (\Psi_r+\ell_q\delta_0^{(r)}),\qquad r\ge1.
\end{equation}
\end{theorem}
\begin{proof}
Fourier coefficients give
$\mathcal S_qD^r=q^rD^r\mathcal S_q$ and
$\mathcal S_q\delta_0^{(r)}=q^{r+1}\delta_0^{(r)}$.
Apply these identities to $F_{m,r}=D^rG_m-a_{m,r}\delta_0^{(r)}$,
use \eqref{dilfp:eq:trace}, and then replace each
$D^rG_j$ by $F_{j,r}+a_{j,r}\delta_0^{(r)}$.
Because $D^r\stconvId=\delta_0^{(r)}$, the contact coefficient is
precisely the displayed $\kappa_{m,r}$. Differentiating the classical
pointwise multiplication row $r$ times proves the raw-coordinate
interpretation of the first sum. Finally $\Psi_r=-F_{0,r}$ and
$\kappa_{0,r}=-\ell_q$, proving \eqref{dilfp:eq:polygamma}.
\end{proof}

These identities answer the normalization question for integer covering
maps and their distribution traces at every index. They also show that
the polynomial convolution algebra is preserved by the normalized trace.
Twisted character projections and products at several independent shifted
singularities still require separate calculations. Exact Fourier-polynomial
and local logarithmic-counterterm controls, with a corruption check, are
recorded in \src{verification/certify_stieltjes_dilation.py}.
